\section{Introdução}\label{sec:intro}

Modelos baseados em Transformers representam a entrada (prompt, documento anexado ou histórico de agente) como uma \textbf{sequência de tokens} processada camada a camada via \emph{self-attention}. Cada token ocupa espaço na janela de contexto e contribui para o custo de inferência, que em atenção densa escala de forma superlinear com o comprimento da sequência\footnote{Em atenção densa, cada token interage com todos os demais; o custo de computação e memória cresce mais rápido que $T$ (tipicamente $O(T^2)$), e não de forma proporcional ($O(T)$).}. Em aplicações com contexto longo (sumarização, RAG, agentes autônomos), boa parte do orçamento de tokens é consumida antes da resposta final: leituras exploratórias, buscas e trechos irrelevantes permanecem no histórico e são reprocessados a cada passo \cite{zhang2026fastcontext,chen2024fastv}.

Evidências recentes quantificam esse gargalo. Zhang \emph{et al.} \cite{zhang2026fastcontext} observam que, em agentes de codificação, leitura e busca consomem cerca de \textbf{46{,}5\%} dos tokens do agente principal; trajetórias end-to-end podem atingir centenas de milhares a milhões de tokens antes de uma correção ser proposta. Separar exploração e devolver contexto compacto reduz o consumo em até \textbf{60\%}, sem degradar (e por vezes melhorando) a taxa de resolução. Benchmarks de \emph{token pruning} em modelos multimodais indicam que muitos tokens carregam pouca informação nova \cite{chen2024fastv,yang2024visionzip}. Comprimir ou selecionar tokens é, portanto, central para custo e qualidade do raciocínio downstream.

Uma linha de trabalho aborda o problema pela \textbf{estrutura algébrica} da matriz de \emph{features} $F \in \mathbb{R}^{T \times D}$ (uma linha por unidade vetorizada, $D$ dimensões por embedding), em vez de heurísticas puramente locais. O SVD-Prune \cite{apedo2026svdprune} critica métodos guiados apenas por scores de atenção: esses pesos tendem a favorecer tokens no início e no fim da sequência (máscara causal, codificação posicional, padrões de leitura), ranqueando tokens menos relevantes quando o orçamento de poda é pequeno. A abordagem decompõe $F$ por SVD, seleciona um subespaço dominante por energia espectral e ranqueia cada token por \emph{leverage scores} das linhas de $U_k$.

Adotamos a SVD \textbf{por documento} (por sequência), e não sobre um corpus agregado: cada texto produz sua matriz $F$; a decomposição identifica modos dominantes \emph{dentro} dessa sequência e os leverage scores ranqueiam tokens \emph{dela}. Isso espelha o uso em inferência: poda de prompt, contexto anexado ou histórico de agente ocorre \textbf{antes} de processar aquela instância, sem acesso a um universo de documentos. Não empilhamos linhas de reviews distintas numa matriz corpus-wide: tal agregação misturaria textos independentes e responderia a outra pergunta (estrutura \textbf{entre} documentos, não poda no eixo sequencial $T \to T'$). A justificativa formal está na Seção~\ref{sec:svd-por-documento}.

Neste relatório de Álgebra Linear Computacional (ALC), isolamos esse operador e comparamos, sob orçamento fixo $\rho$, variantes SVD (\texttt{svd}, \texttt{svd\_energia}) com o baseline \texttt{random}. A pergunta empírica é se seleção por leverage score preserva mais sinal útil do que poda aleatória quando mantemos apenas uma fração $T' \ll T$ das linhas de $F$.

Para quantificar perda de informação, materializamos o fluxo \textbf{texto $\to$ embedding $F$ $\to$ poda $\to$ avaliação} sobre \textbf{reviews longas do IMDb} ($\approx$400 amostras, $\bar{T}\approx 209$ tokens): textos redundantes, sentimento binário supervisionado e comprimento compatível com a motivação de compressão. Cada operador reduz $F$ a $F[S,:]$; treino e teste passam pela mesma poda antes de qualquer medição (Seção~\ref{sec:pipeline}). Usamos dois proxies complementares:
\begin{itemize}
    \item \textbf{Algébrico}: compressão $C$ e, para métodos SVD, energia espectral $E_k^{\mathrm{energia}}$ do subespaço truncado;
    \item \textbf{Supervisionado}: acurácia de sentimento após pooling médio das linhas podadas, normalização $L_2$ e classificador ridge; \texttt{full} (sem poda) serve de teto.
\end{itemize}
Consideramos \emph{menor perda} quando, à compressão equivalente, a acurácia se aproxima mais de \texttt{full} do que nas heurísticas comparadas. Interpretamos o encadeamento teórico (Gram--Schmidt, espectral, $F^\top F$, posto) \cite{alc21svd,alc19espectral,alc11base}; não replicamos agentes nem VLMs completos, e os proxies medem degradação do operador algébrico, não equivalência de saída generativa.

A Seção~\ref{sec:fundamentacao} desenvolve a base teórica; a Seção~\ref{sec:metodologia} detalha o pipeline; a Seção~\ref{sec:resultados} apresenta as comparações empíricas.
